\begin{frame}[fragile]{중요: 목표값의 할인 지수를 세어라}
  \begin{itemize}
    \item $G_t^{(n)}$ 정의에서 $R_{t+k}$의 지수는 $k$. $\tau$번 상태
      갱신 시: \texttt{rewards[tau]} 지수 $0$, \texttt{rewards[tau+1]}
      지수 $1$, …, \texttt{rewards[i]}의 지수는 $\mathbf{i - \tau}$.
    \item 여기서 \textbf{하나만} 틀리면:
  \end{itemize}
  \begin{lstlisting}[basicstyle=\ttfamily\footnotesize]
    G = sum(gamma ** (i - tau - 1) * rewards[i] for i in ...)
    # 각 보상이 gamma^{-1}만큼 "덜" 할인된 -> 편향된 목표값
  \end{lstlisting}
  \begin{itemize}
    \item 이렇게 쓰면 $n=1$일 때조차 \textbf{TD(0)의 목표값이 아님}.
      2000 에피소드 학습 후에도 $V(5) \approx 0.84$로 참값
      $+0.5643$에서 크게 벗어나 \textbf{수렴하지 않음}.
    \item 지수 하나가 틀리면 ``TD(0)를 쓰는 것''이 아니라
      \textbf{``편향된 목표값을 끝없이 쓰는 것''}.
  \end{itemize}
\end{frame}
